\subsection{Algorithm definitions}
\label{app:guidance}
For a free in-grid target command, let $\mathcal N(q)$ be the physical
neighbours of $q$ and $V(q)=|\mathcal N(q)|$. The supplied motion law gives
\[
 \begin{aligned}
 \pi(q)&=\frac{d(q)}{d(q)+\alpha},&
 M(q)&=\pi(q)+(V(q)-1)\frac{1-\pi(q)}{B-1},\\
 \ell(q)&=\frac{\pi(q)}{M(q)},&
 \sigma(q)&=\frac{1-\pi(q)}{(B-1)M(q)}.
 \end{aligned}
\]
\textbf{Value guidance (submitted).} For command $a$ from free cell $p$, the
intended neighbour has weight $\pi(p)$, every other in-grid neighbour
$(1-\pi(p))/(B-1)$, the row is renormalised, and mass on obstacle cells is a
terminal crash. Value iteration uses these exact rows (stored sparsely) with
$r_i=c_{\mathrm{step}}+C\Pr(\text{crash})+G\Pr(\text{arrival at }g_i)$ and
$V_{i,h}(g_i)=0$, because arrival is terminal for that drone. Per-horizon
action values are kept when $nA(H+1)N\le2.5\times10^7$ ($n$ free cells, $A$
commands) on grids of at most 20,000 cells; otherwise the planner uses the
Dijkstra guidance below. On the demo maps this takes under 0.2\,s before
the first decision.

\textbf{Dijkstra guidance (comparison A).} Entry costs are
$1+\lambda\rho(q)$ for a non-goal cell and 1 for the goal, where
\[
 \rho(q)=|\mathcal N(q)\cap O|\sigma(q),\qquad
 \lambda=\frac{w(G+|C|)}
 {\max\{|c_{\mathrm{step}}|,10^{-3}(G+|C|),10^{-9}\}},\quad w=1.
\]
With $w=0$, these are breadth-first distances. The supplied empty-map
distance transform uses a virtual point beyond the grid corner, so slips
remain possible without obstacles.

Grid boundaries are not obstacles: out-of-grid outcome mass is removed
before normalisation. The simulator encodes a joint command as
\[
 a_{\mathrm{int}}=\sum_{i=1}^{N}a_iB^{i-1}.
\]
Maps are built for at most 150,000 cells. Larger grids use Chebyshev
distance without a risk map; fixed-altitude goals on another layer
remain unreachable. Frozen-goal disconnection checks are cached by
drone and reached subset for at most 20,000 cells. Above that size,
the leaf uses static reachability. These fallbacks are implementation
limits, not guarantees of safe coordination.

A rollout cut off after $L$ sampled steps returns
\[
 \sum_{j=0}^{L-1}\gamma^j r_j/\kappa+
 \gamma^L\widehat V(s_L,h-L)/\kappa,
\]
with zero continuation on termination. For an edge sample $v$, update
\[
 n\gets n+1,\quad \Delta\gets v-Q,\quad
 Q\gets Q+\Delta/n,\quad S\gets S+\Delta(v-Q).
\]
The pooled root variance is $\hat\sigma^2=\sum_a S_a/\sum_a(n_a-1)$ over
visited actions, or zero when the denominator is zero.

\subsection{Development-set parameter selection}
\label{app:tuning}
Sixteen layouts cross $N=1,2,3,4$ with
$\alpha\in\{0.25,0.5,1,2\}$ on $10\times10\times1$ or
$10\times10\times5$ grids, with 0--5 obstacles (generator seed 4620).
Each setting uses ten execution seeds and at most 500 simulations per
decision; search-free settings use none. These $\val{DevEpisodes}$
development episodes are separate from all $\val{TotalEpisodes}$ evaluation
episodes.

Round 1 compares risk weights, UCT constants and search-free policies
under most-visited root choice and widening coefficient 1.5.
Round 2 compares the anchored rule, its threshold and rollout exploration,
and widening coefficient 0.5. The selected setting is
\texttt{anchor\_z1\_w0.5}: risk weight 1, UCT constant 0.35,
switch threshold 1, widening coefficient 0.5 and rollout exploration 0.1.
Here the suffix \texttt{w0.5} denotes widening, not the risk weight.
Round 3 compares that Dijkstra guidance with single-drone value-iteration
guidance, each with the anchored rule and with largest-mean-$Q$ root choice,
plus both search-free policies. The selection criterion required no worse
observed success count or mean $L$. Value guidance met it
($\val{TuneVi}$ versus $\val{TuneDijkstra}$ successes), whereas
largest-mean-$Q$ choice raised $L$ under both guidances
($\val{TuneDijkstraMaxq}$ and $\val{TuneViMaxq}$). The round-3 Dijkstra arm
reproduces the round-2 selected setting exactly.
Development tables retain the planner hashes recorded for those versions;
all evaluation suites were rerun with the final implementation.

\begin{table}[H]
\centering\small
\caption{Development results: 160 episodes per setting. $L$: mean
$\pm$95\% interval half-width. $\Delta L$: paired change from the selected
setting, with a 95\% interval over ten seed blocks; lower is better.
Selection across $\val{DevSettings}$ settings makes these exploratory comparisons
optimistic for the development layouts.}
\begin{tabular*}{\linewidth}{@{\extracolsep{\fill}}rlcrr@{}}
\toprule Round & Setting & Success & $L$ & $\Delta L$ (95\% CI) \\
\midrule
\TuningSummaryRows
\bottomrule
\end{tabular*}
\end{table}

\subsection{Factorial results}
\label{app:factorial}
\begingroup
\small
\newcolumntype{R}[1]{>{\raggedleft\arraybackslash}p{#1}}
\newcolumntype{C}[1]{>{\centering\arraybackslash}p{#1}}
\setlength{\LTleft}{0pt}
\setlength{\LTright}{\fill}
\begin{longtable}{@{}R{.055\linewidth}R{.075\linewidth}R{.055\linewidth}
C{.12\linewidth}R{.22\linewidth}R{.22\linewidth}R{.135\linewidth}@{}}
\caption{All 27 S1 conditions, $n=20$ each. $L$ and Return: mean
$\pm$95\% interval half-width; Calls: mean simulator calls per decision.}\\
\toprule $n_y$ & $|A|$ & $|O|$ & Success & $L$ & Return & Calls \\
\midrule
\endfirsthead
\multicolumn{7}{@{}l}{\textit{Table \thetable\ continued}}\\
\toprule $n_y$ & $|A|$ & $|O|$ & Success & $L$ & Return & Calls \\
\midrule
\endhead
\midrule
\multicolumn{7}{r@{}}{\textit{Continued on next page}}\\
\endfoot
\bottomrule
\endlastfoot
\FactorialRows
\end{longtable}
\endgroup

\subsection{Mechanism controls}
\label{app:controls}
Command masks restrict intended directions; the physical slip law retains
all 26 neighbours. Duplicated labels leave the physical commands unchanged.
\begin{table}[H]
\centering\small
\caption{S2 clutter controls ($n_y=10$, $n=20$). $L$: mean $\pm$95\% interval
half-width; Obst./Veh.: collision episodes.}
\begin{tabular*}{\linewidth}{@{\extracolsep{\fill}}llrcrrr@{}}
\toprule Commands & $\alpha$ & $|O|$ & Success & $L$ & Obst. & Veh. \\
\midrule
\ClutterControlRows
\bottomrule
\end{tabular*}
\end{table}
\begin{table}[H]
\centering\small
\caption{Duplicated labels ($n_y=10$, $|O|=2$, $n=20$).
Root: admitted actions; Selected: visit share of the executed action.
$L$: mean $\pm$95\% interval half-width.}
\begin{tabular*}{\linewidth}{@{\extracolsep{\fill}}lrcrrr@{}}
\toprule Labels & $|A|$ & Success & $L$ & Root & Selected (\%) \\
\midrule
\BranchControlRows
\bottomrule
\end{tabular*}
\end{table}

\subsection{Component effects}
\label{app:components}
\begin{table}[H]
\centering\small
\caption{Overall paired S3 effects, with 95\% intervals.
Uniform proposals: ten seeds; others: twenty.}
\begin{tabular*}{\linewidth}{@{\extracolsep{\fill}}lll@{}}
\toprule Variant & $\Delta L$ & $\Delta$success (pp) \\
\midrule
\ComponentOverallRows
\bottomrule
\end{tabular*}
\end{table}
\begin{table}[H]
\centering\small
\caption{Mean paired $\Delta L$ at low/high complexity levels.
$^\ast$: the interaction's 95\% interval excludes zero.}
\begin{tabular*}{\linewidth}{@{\extracolsep{\fill}}lccc@{}}
\toprule Variant & $n_y$: 5 / 20 & $|A|$: 36 / 676 & $|O|$: 0 / 5 \\
\midrule
\ComponentComplexityRows
\bottomrule
\end{tabular*}
\end{table}

\begin{table}[H]
\centering\small
\setlength{\tabcolsep}{3pt}
\caption{S3 successes at all eight factorial corners. Uniform proposals:
$n=10$ per corner; all other rows: $n=20$. Within each $n_y$ group, columns
give intended commands per drone / obstacle count. These counts complement
the paired effects and interaction intervals above.}
\begin{tabular*}{\linewidth}{@{\extracolsep{\fill}}lcccccccc@{}}
\toprule
& \multicolumn{4}{c}{$n_y=5$} & \multicolumn{4}{c}{$n_y=20$}\\
\cmidrule(lr){2-5}\cmidrule(l){6-9}
Variant & 6/0 & 6/5 & 26/0 & 26/5 & 6/0 & 6/5 & 26/0 & 26/5 \\
\midrule
\ComponentRawRows
\bottomrule
\end{tabular*}
\end{table}

\subsection{Search-budget diagnostics}
\label{app:budget}
Root branching and mean new-node depth measure how the search changes as
the simulation cap grows. New-node depth is a sampled mean, not the maximum
configured tree depth of eight. The four-drone tree remains shallow even
at 2,000 trials, consistent with limited reuse of stochastic joint successors.
\begin{table}[H]
\centering\small
\caption{S4 diagnostics, $n=20$ per row. $L$: mean $\pm$95\% interval
half-width; Root: mean admitted actions; Depth: mean new-node depth.
Drone counts are in parentheses.}
\begin{tabular*}{\linewidth}{@{\extracolsep{\fill}}lrcrrr@{}}
\toprule Task (drones) & Sims & Success & $L$ & Root & Depth \\
\midrule
\BudgetRows
\bottomrule
\end{tabular*}
\end{table}

\begin{samepage}
\subsection{Regression coverage and paired baseline}
\label{app:regression}
The 18 named scenarios each use 20 seeds and a cap of 2,000 simulations.
They cover initial goal occupancy, swaps, boundaries, small or thin grids,
unreachable goals, a short horizon, narrow passages, altitude changes,
strong noise and multi-drone motion. The three statically infeasible cases
and the one-step horizon are expected failures; collisions on feasible
tasks remain recorded failures.
\par
\end{samepage}

\begingroup
\small
\setlength{\LTleft}{0pt}
\setlength{\LTright}{\fill}
\begin{longtable}{@{}>{\raggedright\arraybackslash}p{.28\linewidth}
>{\centering\arraybackslash}p{.09\linewidth}
>{\centering\arraybackslash}p{.11\linewidth}
>{\raggedleft\arraybackslash}p{.055\linewidth}
>{\raggedleft\arraybackslash}p{.055\linewidth}
>{\raggedleft\arraybackslash}p{.165\linewidth}
>{\raggedleft\arraybackslash}p{.12\linewidth}@{}}
\caption{Named regression scenarios. No search: paired base policy;
Obst./Veh.: collision episodes. $L$: mean $\pm$95\% interval half-width,
using each scenario's own horizon as the failure penalty. Calls/step
counts all generative-model calls, including rollouts.}\\
\toprule Scenario & Success & No search & Obst. & Veh. & $L$ & Calls/step \\
\midrule
\endfirsthead
\multicolumn{7}{@{}l}{\textit{Table \thetable\ continued}}\\
\toprule Scenario & Success & No search & Obst. & Veh. & $L$ & Calls/step \\
\midrule
\endhead
\midrule
\multicolumn{7}{r@{}}{\textit{Continued on next page}}\\
\endfoot
\bottomrule
\endlastfoot
\RegressionRows
\end{longtable}
\endgroup

Among $\val{BaseMultiPairs}$ feasible multi-drone regression/generated pairs,
search succeeds in $\val{BaseMultiFull}$ versus $\val{BaseMultiGreedy}$;
removing search changes $L$ by $\val{BaseMultiDL}$ $\val{BaseMultiDLci}$.
Search rescues $\val{BaseMultiRescued}$ episodes and loses
$\val{BaseMultiLost}$. On the $\val{BaseSinglePairs}$
single-drone pairs, search succeeds in $\val{BaseSingleFull}$ versus
$\val{BaseSingleGreedy}$ ($\Delta L=\val{BaseSingleDL}$ $\val{BaseSingleDLci}$).
These exploratory intervals resample whole execution-seed clusters,
retaining the episode-weighted means across layouts; they are distinct from
the balanced S1--S3 seed-block $t$ intervals. Removing search changes
multi-drone success by $\val{BaseMultiDP}$\,pp $\val{BaseMultiDPci}$.
Both the multi-drone success and completion-length intervals favour the
base policy on these fixed layouts, without establishing a general ordering.

An exhaustive small-model check separately verifies 704 state-action
pairs and 6,048 positive-probability outcomes on all 11 admissible
nonterminal states of a free $2\times2\times1$ two-drone grid
(\texttt{results/small\_model\_verification.json}). This checks model
semantics on that finite example, not performance on arbitrary maps.

\subsection{Value-iteration versus Dijkstra guidance}
\label{app:guidance-compare}
S1, S2, S4, S5 and the regression/supplementary suites compare the earlier
Dijkstra guidance (A) and submitted value guidance (B) on the same scenarios
and seeds; pairs match
scenario, variant, simulation cap and seed. S3 is omitted here because
removing risk terms has different meanings under A and B. The A runs are archived in
\texttt{results/guidance\_dijkstra/}. Over all solver suites, B succeeds in
$\val{GuideB}$ of $\val{GuidePairs}$ distinct paired episodes and A in
$\val{GuideA}$: B rescues $\val{GuideWins}$ and loses $\val{GuideLosses}$.
The success change is $\val{GuideDP}$\,pp $\val{GuideDPci}$; penalised length
changes by $\val{GuideDL}$ $\val{GuideDLci}$ steps. The aggregate excludes
$\val{GuideDuplicates}$ repeated conditions shared by the budget and other
suites. Intervals use 20,000 execution-seed cluster bootstrap resamples,
rather than treating all episode differences as independent. Value guidance
was selected on the development set before these evaluation runs; the
comparison does not establish uniform superiority over Dijkstra guidance.
The pooled completion-length interval favours value guidance, while the
success-change interval includes zero.

\begin{table}[H]
\centering\small
\setlength{\tabcolsep}{3pt}
\caption{Paired comparison of guidance B (value iteration, submitted) with A
(Dijkstra). Rescues/Losses: pairs that only B / only A completes; $\Delta L$:
B minus A with a 95\% seed-cluster bootstrap interval.
The aggregate counts repeated conditions once.}
\begin{tabular*}{\linewidth}{@{\extracolsep{\fill}}lrrrrrl@{}}
\toprule Suite & Pairs & A & B & Rescues & Losses & $\Delta L$ (95\% CI) \\
\midrule
\GuidanceRows
\bottomrule
\end{tabular*}
\end{table}

\subsection{Supplementary end-to-end cases}
\label{app:supplement}
The supplementary cases below include fixed-altitude single-drone maps,
separate altitude layers and a goal inside a teammate's corridor.
Wall-gap and corridor cases use seven/eight obstacles, beyond the demo's
five-obstacle range. No timeout occurred; all failures were collisions.
The DP maximum is the largest success probability any policy can achieve
within $H$ under the supplied model, computed by finite-horizon value iteration
(\texttt{scripts/compute\_ceilings.py}; the two-layer crossing case is a product
because drones on different fixed layers never interact; the four-drone case
was not solved). This optimises success probability, whereas B1 optimises
discounted reward. The computed maxima lie inside the solver's Wilson
intervals, which does not establish policy equivalence.
\begin{table}[H]
\centering\small
\setlength{\tabcolsep}{3pt}
\caption{Supplementary cases ($n=20$ per policy). 95\% CI: Wilson interval of
the success rate; No search: paired base
policy; Obst./Veh./T.o.: full-solver failure counts.
$L$: mean $\pm$95\% interval half-width.}
\begin{tabular*}{\linewidth}{@{\extracolsep{\fill}}lcccrrrr@{}}
\toprule Case & Success & 95\% CI (\%) & No search & Obst. & Veh. & T.o. & $L$ \\
\midrule
\SupplementRows
\bottomrule
\end{tabular*}
\end{table}

\begin{table}[H]
\centering\small
\caption{Enumerated finite-horizon success maxima (floating-point
arithmetic) and policy validation in the supplied simulator.
Five policies use seeds 0--1,999; the corridor adds a fixed batch of
8,000 disjoint seeds, giving 10,000 runs in total. The initial corridor
interval missed its DP value; the original batch is retained in
\texttt{results/ceilings.json} and the extension in
\texttt{results/ceiling\_validation.json}. Wilson 95\% intervals describe
these diagnostic samples, separate from the $\val{TotalEpisodes}$ solver-evaluation episodes;
they do not establish policy equivalence.}
\begin{tabular*}{\linewidth}{@{\extracolsep{\fill}}lrrr@{}}
\toprule Case & DP max. (\%) & Policy successes & 95\% CI (\%) \\
\midrule
\CeilingRows
\bottomrule
\end{tabular*}
\end{table}

\subsection{Generated-layout results}
\label{app:layouts}
\begin{table}[H]
\centering\small
\caption{All 16 S5 layouts, one per $(N,\alpha)$ pair, ten seeds each.
$L$: mean $\pm$95\% interval half-width; Depth: mean new-node depth.}
\begin{tabular*}{\linewidth}{@{\extracolsep{\fill}}rrcrrr@{}}
\toprule $N$ & $\alpha$ & Success & Wilson 95\% (\%) & $L$ & Depth \\
\midrule
\RandomCaseRows
\bottomrule
\end{tabular*}
\end{table}
